\documentclass[a4paper,11pt]{article}
\usepackage[T1]{fontenc}
\usepackage[utf8]{inputenc}
\usepackage[brazilian]{babel}
\usepackage{lmodern}
\usepackage{textcomp}
\usepackage[margin=1.8cm,top=2.0cm,bottom=2.0cm,headheight=15pt]{geometry}
\usepackage{xcolor,graphicx,booktabs,tabularx,array,enumitem,listings,fancyhdr,tikz}
\usepackage[unicode,colorlinks=true,urlcolor=teal,linkcolor=teal,breaklinks=true]{hyperref}
\usetikzlibrary{arrows.meta,positioning}
\definecolor{azul}{HTML}{12334B}
\definecolor{verde}{HTML}{007E87}
\definecolor{suave}{HTML}{EDF3F7}
\definecolor{cinza}{HTML}{455D6C}
\renewcommand{\familydefault}{\sfdefault}
\setlength{\parindent}{0pt}
\setlength{\parskip}{6pt}
\setlength{\emergencystretch}{2em}
\setlist{nosep,leftmargin=1.4em,itemsep=3pt,topsep=4pt}
\renewcommand{\arraystretch}{1.25}
\newcolumntype{Y}{>{\raggedright\arraybackslash}X}
\pagestyle{fancy}
\fancyhf{}
\fancyhead[L]{\small\color{cinza}IFES / REDES PARA CONTROLE E AUTOMAÇÃO}
\fancyhead[R]{\small\color{verde}2026/2}
\fancyfoot[L]{\small\color{cinza}\docshort}
\fancyfoot[R]{\small\thepage\ / \pageref{fim}}
\renewcommand{\headrulewidth}{0.3pt}
\lstset{basicstyle=\ttfamily\small,backgroundcolor=\color{suave},
 frame=single,rulecolor=\color{suave},framesep=6pt,columns=fullflexible,
 keepspaces=true,showstringspaces=false,upquote=true,breaklines=true,breakatwhitespace=false,
 aboveskip=7pt,belowskip=7pt,tabsize=2,
 literate={á}{{\'a}}1 {à}{{\`a}}1 {ã}{{\~a}}1 {â}{{\^a}}1
 {é}{{\'e}}1 {ê}{{\^e}}1 {í}{{\'i}}1 {ó}{{\'o}}1 {ô}{{\^o}}1
 {õ}{{\~o}}1 {ú}{{\'u}}1 {ç}{{\c c}}1 {Á}{{\'A}}1 {É}{{\'E}}1
 {Ç}{{\c C}}1 {Ã}{{\~A}}1 {→}{{$\to$}}1}
\newcommand{\capa}[3]{
 {\small\color{verde}\bfseries LABORATÓRIO DE REDES\hfill #1}\par\vspace{0.6cm}
 {\Huge\color{azul}\bfseries #2\par}\vspace{0.25cm}
 {\Large\color{cinza}#3\par}\vspace{0.5cm}}
\newcommand{\pagina}[2]{\newpage{\small\color{verde}\bfseries #1}\par
 {\LARGE\color{azul}\bfseries #2\par}\vspace{0.2cm}}
\newcommand{\sub}[1]{\par\vspace{0.15cm}{\large\color{verde}\bfseries #1}\par}
\newcommand{\caixa}[2]{\par\vspace{4pt}\noindent\colorbox{suave}{\parbox{\dimexpr\linewidth-2\fboxsep\relax}{\textbf{\color{azul}#1}\par #2}}\par\vspace{4pt}}
\newcommand{\fonte}[2]{{\footnotesize\color{cinza}Fonte: \href{#1}{#2}.}\par}
\newcommand{\cmd}[1]{\texttt{\detokenize{#1}}}
\newcommand{\finalizar}{\label{fim}\end{document}}

\newcommand{\docshort}{Prática 2 · API REST, Postman e interface web}
\hypersetup{pdftitle={Prática 2 — CRUD REST com Postman e interface web},pdfauthor={Material didático — Redes / Ifes Campus Serra}}
\begin{document}
\capa{PRÁTICA 02}{Do pedido ao cadastro}{Consumir uma API REST com Postman e uma interface web}
\textbf{Curso:} Engenharia de Controle e Automação · Ifes, Campus Serra\\
\textbf{Duração sugerida:} 120 minutos de investigação REST, após iniciar o kit; reserve 20--30 minutos para a primeira instalação na EC2.\\
\textbf{Organização:} duplas/trios, com uma EC2 por equipe e credenciais geradas automaticamente.

\sub{O desafio}
Você iniciará uma aplicação pronta de inventário de equipamentos fictícios das Plantas A, B e C. Deverá cadastrar, consultar, alterar e excluir recursos; depois, comparar as requisições do Postman com as enviadas por uma interface web.

\caixa{Aplicação pronta, foco em REST}{O repositório contém API, banco, página web, coleção Postman e scripts prontos. Você apenas inicia o kit na sua EC2 e valida o serviço. Não é necessário programar, gerar interface com IA nem receber uma chave do professor. A mesma API será consumida pelo Postman e pelo navegador.}

\sub{Ao terminar, você deverá conseguir}
\begin{itemize}
\item Relacionar operações CRUD a métodos, caminhos, corpos e status HTTP.
\item Enviar JSON e credenciais no lugar correto; interpretar erros.
\item Distinguir substituição completa por PUT de alteração parcial por PATCH.
\item Comprovar persistência e isolamento entre grupos.
\item Explicar a função do proxy reverso, de HTTPS e de CORS.
\end{itemize}

\sub{Antes de iniciar}
Tenha acesso a uma conta AWS autorizada, uma EC2 Ubuntu Server 24.04 LTS dedicada, sua chave SSH e o Postman Desktop na estação. A EC2 pode ser a da prática 1; o kit também funciona em uma instância nova. A conta precisa permitir EC2 e acesso à Internet para baixar o repositório e os pacotes.

\begin{tabularx}{\linewidth}{@{}lY@{}}\toprule
Você obtém & Onde encontrar \\\midrule
Repositório e kit & \href{https://github.com/rafaelrezo/pratica-rest-ec2}{github.com/rafaelrezo/pratica-rest-ec2}.\\
DNS e IPv4 públicos & Console AWS, nos detalhes da sua EC2.\\
Chaves da API & Geradas pelo script na própria EC2; nunca vêm do professor.\\
Certificado público & Gerado pelo script e copiado por SSH verificado.\\\bottomrule
\end{tabularx}
\textbf{Roteiro:} iniciar o kit; validar o acesso; explorar o contrato no Postman; usar a página pronta; comparar as mensagens HTTP. Tudo opera sobre dados fictícios: mudar \cmd{ativo} não liga ou desliga uma máquina.

\pagina{INICIAR / KIT PRONTO}{Do repositório à aplicação na EC2}
\textbf{No console AWS:} use Ubuntu Server 24.04 LTS em uma instância dedicada, com IPv4 público, DNS público e sub-rede com rota para Internet Gateway. Libere TCP 22 somente para seu IP/rede administrativa e TCP 80/443 para o acesso web; não abra 8000. Se o UFW já estiver ativo, permita OpenSSH e Nginx Full. Confira custos, permissões e o término dos recursos da equipe.

\sub{1. Conectar e obter os artefatos}
Na estação, substitua os marcadores e confirme a chave do host SSH pelo console EC2 antes de aceitar a primeira conexão:
\begin{lstlisting}
ssh -i SUA_CHAVE.pem ubuntu@SEU_DNS_EC2
\end{lstlisting}
Na EC2, obtenha o repositório pronto:
\begin{lstlisting}
sudo apt update
sudo apt install -y git
git clone https://github.com/rafaelrezo/pratica-rest-ec2.git pratica-rest-ec2
cd ~/pratica-rest-ec2
\end{lstlisting}
Sem Git, envie a pasta \cmd{kit/} extraída do ZIP por SCP para \cmd{~/pratica-rest-ec2/kit/}. Os scripts são os mesmos; não copie trechos de backend do PDF.

\sub{2. Iniciar e validar}
Ainda na EC2, execute os dois comandos. O início identifica automaticamente o DNS e o IPv4 públicos pelos metadados da instância:
\begin{lstlisting}
bash kit/iniciar-lab.sh
bash kit/validar-lab.sh
\end{lstlisting}
O primeiro script prepara os serviços, o banco, HTTPS e duas identidades de teste. O segundo mostra método, caminho e status esperado/recebido, criando e removendo apenas seu item de ensaio. \textbf{Esperado:} saúde 200, sem chave 401, criação 201, leitura/alteração 200, PUT incompleto 400, isolamento 404, exclusão 204 e consulta final 404.

\caixa{Validar não encerra a prática}{O script confirma que o laboratório está pronto. Nas etapas seguintes, você fará as requisições no Postman e na página para explicar o que cada mensagem significa. Não precisa editar Python, HTML ou configuração do Nginx.}

\textbf{Se repetir a inicialização:} as chaves e o banco são mantidos; a página de referência é republicada. Se o DNS/IP mudar ou o certificado expirar, execute novamente e copie o certificado atualizado. Se a descoberta automática estiver desabilitada, use \cmd{bash kit/iniciar-lab.sh SEU_DNS_EC2 IP_PUBLICO} com os valores atuais do console. A instalação usa Internet e pode levar alguns minutos.

\pagina{PREPARAR / ACESSO}{Obter sua chave e conferir HTTPS}
\sub{1. Consultar as credenciais locais}
Na sessão SSH privada, fora de gravações/capturas de tela:
\begin{lstlisting}
cat ~/chaves-redes/entrega-grupos.csv
openssl x509 -in ~/redes.crt -noout -fingerprint -sha256
\end{lstlisting}
O CSV contém \cmd{g01} e \cmd{g02}, com chaves aleatórias distintas. Use \textbf{g01 no CRUD principal}; g02 é uma segunda identidade para testar isolamento na mesma API. Não são duas equipes físicas e não é preciso combinar o teste com colegas. Cada EC2 gera suas próprias chaves, mesmo repetindo esses nomes.

\textbf{Três arquivos diferentes:} \cmd{SUA_CHAVE.pem} permite SSH; \cmd{redes.crt} é o certificado público HTTPS; a chave de g01 autentica requisições em \cmd{X-API-Key}. A chave privada TLS fica no servidor. Não as confunda.

\sub{2. Obter somente o certificado público}
Em outro terminal, \textbf{na estação}:
\begin{lstlisting}
scp -i SUA_CHAVE.pem ubuntu@SEU_DNS_EC2:~/redes.crt .
openssl x509 -in redes.crt -noout -fingerprint -sha256
curl --cacert redes.crt -i https://SEU_DNS_EC2/api/saude
\end{lstlisting}
Compare as impressões digitais da estação e da sessão SSH. Espere 200 e \cmd{{"status":"ok"}}. O certificado autoassinado tem validade de sete dias e cobre o DNS/IPv4 informados. O DNS da AWS não fornece confiança HTTPS automaticamente.

\sub{3. Abrir a aplicação pronta}
Abra \cmd{https://SEU_DNS_EC2/painel.html}. Confira o certificado e estabeleça a confiança específica no navegador. Digite a chave de g01 e clique em ``Conectar / listar''. A página deve mostrar uma lista vazia ou os registros já criados nessa identidade. A chave fica apenas na memória da página.

\textbf{Se o acesso externo falhar:} o script já testou aplicação e HTTPS dentro da EC2. Confira DNS/IP atuais, regras 443, rota pública e rede da estação. Erro de certificado exige conferir SAN, data e confiança. Não desligue a verificação TLS globalmente nem exponha 8000.

\caixa{Ponto de partida}{Anote a URL, o GET de saúde e o resultado da validação automática, sem registrar chaves. O restante da prática investiga a API que já está funcionando.}

\pagina{ENTENDER A ARQUITETURA}{Uma API, vários clientes}
\begin{center}
\begin{tikzpicture}[node distance=14mm,box/.style={draw=verde,rounded corners,fill=suave,align=center,text width=3.7cm,minimum height=1.6cm},>={Stealth}]
\node[box](p){Postman ou navegador\\cliente HTTP};
\node[box,right=of p](n){Nginx na EC2\\porta pública 443};
\node[box,right=of n](a){Gunicorn + Flask\\127.0.0.1:8000};
\draw[->,thick](p)--node[above,font=\small]{HTTPS}(n);
\draw[->,thick](n)--node[above,font=\small]{HTTP}(a);
\node[below=5mm of a,font=\small,align=center]{SQLite em disco\\cadastros por grupo};
\end{tikzpicture}
\end{center}
O Nginx entrega a página em \cmd{/} e encaminha \cmd{/api/} para a aplicação. Gunicorn executa o aplicativo Flask como um serviço. SQLite persiste os dados em arquivo. A porta 8000 escuta apenas em loopback e não é liberada no Security Group.

\sub{O que torna esta interface orientada a recursos?}
\cmd{/api/equipamentos} identifica uma coleção. \cmd{/api/equipamentos/7} identifica um item. Os caminhos usam substantivos e os métodos expressam ações. As representações são JSON. As respostas utilizam status HTTP e cada requisição protegida leva a credencial necessária.

Este é um CRUD didático inspirado no estilo REST. Usar JSON e métodos HTTP, isoladamente, não demonstra todas as restrições de uma arquitetura REST completa. O objetivo da aula é aplicar semântica HTTP, recursos e representações de forma consistente.

\sub{Como o grupo é identificado?}
O cabeçalho \cmd{X-API-Key} contém a chave de g01 gerada na sua EC2. O servidor associa essa chave ao grupo e restringe todas as operações aos seus registros. Não existe campo \cmd{grupo} no corpo; tentar editar um item de outro grupo retorna 404.

\textbf{Limites:} a chave identifica o grupo, não uma pessoa. Trata-se de autenticação didática com autorização por grupo, sem gestão individual de usuários. Ela é enviada por HTTPS; a interface do aluno não consegue esconder de seu próprio usuário a chave que utiliza. Não a publique no código nem em prompts.

\caixa{Pense antes do primeiro teste}{Quem inicia a conexão? Qual programa termina TLS? Quem lê e grava SQLite? Por que CORS não substitui a chave? Identifique também onde os dados ainda podem aparecer em texto claro: cliente, Nginx, loopback e processo da aplicação.}
\fonte{https://flask.palletsprojects.com/en/stable/deploying/gunicorn/}{Flask: execução com Gunicorn}

\pagina{CONTRATO / DADOS}{Conheça o recurso antes de enviar}
\begin{lstlisting}
{
  "id": 7,
  "nome": "CLP-A-01",
  "tipo": "clp",
  "local": "Planta A",
  "ativo": true
}
\end{lstlisting}
\begin{tabularx}{\linewidth}{@{}lY@{}}\toprule
Campo & Regra desta API \\\midrule
\cmd{id} & Gerado pelo servidor. Usar na URL; não enviar no corpo.\\
\cmd{nome} & Texto com 1--80 caracteres após retirar espaços das pontas; único por grupo.\\
\cmd{tipo} & Um de: \cmd{sensor}, \cmd{clp}, \cmd{ihm}, \cmd{atuador}, \cmd{outro}.\\
\cmd{local} & Texto com 1--80 caracteres. Ex.: Planta A, Operação, IDMZ.\\
\cmd{ativo} & Booleano JSON: \cmd{true} ou \cmd{false}, sem aspas.\\\bottomrule
\end{tabularx}

\sub{Rotas disponíveis}
\begin{tabularx}{\linewidth}{@{}llY@{}}\toprule
Método & Caminho após \cmd{/api} & Resposta de sucesso \\\midrule
GET & \cmd{/saude} & 200 \cmd{{"status":"ok"}}, sem chave.\\
GET & \cmd{/equipamentos} & 200 com \cmd{itens} e \cmd{total}.\\
GET & \cmd{/equipamentos/{id}} & 200 com um objeto.\\
POST & \cmd{/equipamentos} & 201 com objeto e \cmd{Location}.\\
PUT & \cmd{/equipamentos/{id}} & 200 com objeto substituído.\\
PATCH & \cmd{/equipamentos/{id}} & 200 com objeto atualizado.\\
DELETE & \cmd{/equipamentos/{id}} & 204, sem corpo.\\\bottomrule
\end{tabularx}
Todas as rotas, exceto saúde e preflight OPTIONS, exigem chave. Use caminhos sem barra final. Não há filtros, paginação nem busca no servidor. POST e PUT exigem os quatro campos editáveis; PATCH aceita um ou mais. Campos desconhecidos são rejeitados. O limite de corpo é 16 KiB.

\textbf{Exemplo de listagem vazia:} \cmd{{"itens":[],"total":0}}. O ID 7 é apenas ilustrativo. Não presuma que o primeiro cadastro receberá ID 1: os IDs são gerados pela API e compartilhados pelas identidades g01/g02 dessa instância.

\pagina{POSTMAN / PREPARAR}{Configurar ambiente e confiança TLS}
\begin{enumerate}
\item Abra o Postman Desktop. Pelo botão \textit{Import}, carregue a coleção:\newline
\cmd{redes-crud.postman_collection.json}; depois, o ambiente:\newline
\cmd{redes.postman_environment.json}. Ambos estão em \cmd{kit/postman/}.
\item Selecione o ambiente ``Redes | preencher localmente''. Defina \cmd{baseUrl} com a URL HTTPS da sua EC2, \textbf{sem} \cmd{/api} e sem barra final.
\item Defina \cmd{apiKey} com a chave de g01 do seu CSV, como valor local/privado ou em um Vault, conforme a versão disponível. Marcar como segredo oculta a exibição, mas não assegura, sozinho, que o valor nunca seja sincronizado. Não compartilhe ambiente preenchido.
\item Deixe \cmd{equipamentoId} vazio; a requisição de criação vai preenchê-lo usando o ID da resposta.
\item Se houver certificado público válido, mantenha a verificação TLS padrão. Se autoassinado, obtenha \cmd{redes.crt} por canal verificado e adicione-o em \textit{Settings → Certificates → CA Certificates} no aplicativo Desktop. Verifique SAN e validade. Mantenha \textit{SSL certificate verification} ligada.
\end{enumerate}
\fonte{https://learning.postman.com/docs/getting-started/basics/about-postman-agent/}{Postman: Desktop e agentes}

\sub{Faça o primeiro GET manualmente}
Crie uma requisição, selecione GET e informe:
\begin{lstlisting}
{{baseUrl}}/api/saude
\end{lstlisting}
Clique \textit{Send}. A resposta deve ser 200, com \cmd{{"status":"ok"}}. Esse teste verifica que Nginx e aplicação respondem pela rede, mas ainda não testa sua chave nem uma operação de banco.

\sub{Prepare a autenticação das chamadas protegidas}
Na coleção importada, a aba \textit{Authorization} já usa \textit{API Key}: chave \cmd{X-API-Key}, valor \cmd{{{apiKey}}}, local \textit{Header}. Requisições herdam essa configuração; os testes sem chave a desativam de propósito.

Para uma requisição nova fora da coleção, configure manualmente o cabeçalho \cmd{X-API-Key: SUA_CHAVE}. Nunca adicione a chave à query string. Ao listar \cmd{/api/equipamentos}, o resultado deve ser 200. Sem chave, deve ser 401.

\caixa{Verificação 1}{Registre: sucesso de saúde; tentativa sem chave; listagem autenticada. O relatório deve ocultar o valor de \cmd{X-API-Key}. Se saúde responde 200 e listar retorna 401, explique qual parte da infraestrutura já demonstrou funcionar.}

\pagina{CRUD / CREATE E READ}{Criar e conferir um cadastro}
\sub{1. POST para a coleção}
No Postman: método POST; URL \cmd{{{baseUrl}}/api/equipamentos}; aba \textit{Body → raw → JSON}. O cabeçalho \cmd{Content-Type: application/json} deve estar presente.
\begin{lstlisting}
{
  "nome": "CLP-A-01",
  "tipo": "clp",
  "local": "Planta A",
  "ativo": true
}
\end{lstlisting}
\textbf{Esperado:} 201 Created; objeto com ID gerado; cabeçalho \cmd{Location: /api/equipamentos/ID}. Se o nome já existir em seu grupo, altere-o ou use um cadastro que você ainda não criou.

\sub{2. GET para o item criado}
Copie o ID real retornado para a variável \cmd{equipamentoId}. Na coleção automática, o teste do POST faz isso por você.
\begin{lstlisting}
GET {{baseUrl}}/api/equipamentos/{{equipamentoId}}
X-API-Key: {{apiKey}}
\end{lstlisting}
O resultado deve ser 200 com os dados enviados e o ID. Faça também GET na coleção e confirme que \cmd{itens} contém o cadastro e que \cmd{total} corresponde à quantidade retornada.

\sub{3. A mesma ideia com curl, opcional}
Para comparar ferramentas, configure na estação uma variável temporária, sem colocar a chave no histórico. Use um shell Bash compatível:
\begin{lstlisting}
BASE_URL='https://SEU_DNS_EC2'
read -rsp 'Chave do grupo: ' API_KEY; echo
curl --cacert redes.crt -i "$BASE_URL/api/equipamentos" \
  -H "X-API-Key: $API_KEY"
unset API_KEY
# redes.crt foi obtido pelo canal verificado da pratica 1.
\end{lstlisting}
O cabeçalho será enviado ao servidor; mantenha o terminal privado. O \cmd{curl} e o Postman podem representar a mesma requisição. Não é necessário enviar pelo navegador para usar HTTP.

\caixa{Verificação 2}{Mostre que criar e ler são requisições diferentes. Identifique por que o POST recebe JSON no corpo e por que o ID é obtido da resposta. Explique por que repetir o POST pode produzir 409, enquanto repetir o GET não cria outro equipamento.}

\pagina{CRUD / UPDATE}{PUT e PATCH não pedem a mesma coisa}
\sub{1. PATCH: alterar apenas um campo}
Método PATCH na URL do item. Corpo JSON:
\begin{lstlisting}
{"ativo": false}
\end{lstlisting}
\textbf{Esperado:} 200; \cmd{ativo} fica falso; nome, tipo e local permanecem. Leia novamente por GET para conferir. O campo apenas informa o estado do cadastro.

\sub{2. PUT: enviar toda a representação editável}
Método PUT na mesma URL. Corpo JSON:
\begin{lstlisting}
{
  "nome": "CLP-A-01",
  "tipo": "clp",
  "local": "Planta B",
  "ativo": true
}
\end{lstlisting}
\textbf{Esperado:} 200, com todos os campos fornecidos. Neste contrato, PUT substitui um item existente; não cria um ID desconhecido. Se o item não existir, retorna 404. Esse comportamento está definido pela API e não deve ser adivinhado pelo cliente.

\sub{3. Teste a diferença de propósito}
Envie somente \cmd{{"ativo":false}} por PUT. A API deve responder 400, pois faltam campos obrigatórios. Envie o mesmo corpo por PATCH: deve responder 200. Não generalize que todo backend valida de forma idêntica; o cliente deve seguir o contrato recebido.

\sub{4. Faça uma previsão antes de executar}
\begin{tabularx}{\linewidth}{@{}lYY@{}}\toprule
Operação & Previsão de efeito & Evidência após o GET \\\midrule
PATCH \cmd{ativo=false} & \rule{2.8cm}{0.2pt} & \rule{3cm}{0.2pt}\\
Mesmo PATCH novamente & \rule{2.8cm}{0.2pt} & \rule{3cm}{0.2pt}\\
PUT completo & \rule{2.8cm}{0.2pt} & \rule{3cm}{0.2pt}\\
PUT incompleto & \rule{2.8cm}{0.2pt} & \rule{3cm}{0.2pt}\\\bottomrule
\end{tabularx}

\caixa{Verificação 3}{Justifique por que atribuir o mesmo valor com PATCH pode ser idempotente. Explique também por que o corpo de PUT inclui \cmd{ativo}, mesmo quando você só pretendia alterar \cmd{local}.}

\pagina{CRUD / DELETE E ERROS}{Resultados esperados também incluem falhas}
\sub{1. DELETE do cadastro criado pelo grupo}
Envie DELETE para \cmd{/api/equipamentos/ID}. \textbf{Esperado:} 204 No Content, sem corpo. Depois, GET para o mesmo ID deve retornar 404. Repetir DELETE também retorna 404; o recurso continua ausente.

\sub{2. Execute estes testes controlados}
\begin{tabularx}{\linewidth}{@{}p{4.7cm}lY@{}}\toprule
Alteração no teste & Status & Explicação esperada \\\midrule
Retirar a chave de uma rota protegida & 401 & Não foi apresentada credencial válida.\\
POST incompleto ou JSON inválido & 400 & A representação recebida é inválida.\\
Enviar \cmd{ativo: "true"} no JSON & 400 & Texto não é booleano.\\
Repetir nome no mesmo grupo & 409 & Conflito com recurso existente.\\
Enviar \cmd{text/plain} no POST & 415 & Formato de corpo não aceito.\\
PUT na coleção, sem ID & 405 & Método não disponível nesse recurso.\\
GET de item excluído & 404 & Esse item não está disponível no grupo.\\\bottomrule
\end{tabularx}
Erros da aplicação têm corpo \cmd{{"erro":"mensagem"}}. Erros do Nginx, como 502, podem vir em HTML. Por isso, o cliente precisa verificar status e conteúdo; não deve presumir que toda resposta sempre é JSON.

\sub{3. Isolamento sem depender de outra equipe}
Crie um novo item com a chave de g01 e anote o ID. Duplique o ambiente Postman e use nele a chave de g02 do seu CSV, mantendo \cmd{baseUrl}. Para o ID de g01, GET, PATCH, PUT e DELETE com g02 devem retornar 404. Volte a g01: GET deve retornar 200 e os dados devem estar intactos. Exclua o item com g01. A segunda chave é válida; o item é que está fora de seu escopo.

\caixa{Verificação 4}{Preencha uma tabela com método, caminho, status, corpo resumido e interpretação para criar, listar, consultar, alterar, excluir e dois erros. Não inclua chaves. Uma resposta 404 é evidência de que o servidor respondeu HTTP; ela não equivale a timeout de rede.}

\textbf{Limite de segurança do exercício:} a API possui validação e isolamento por grupo, mas não foi projetada para uso de produção. Não utilize dados pessoais nem a conecte a CLPs, sensores ou atuadores reais.

\pagina{POSTMAN / REPETIR}{Executar a coleção e ler os testes}
\sub{Sequência da coleção fornecida}
\begin{tabularx}{\linewidth}{@{}llY@{}}\toprule
Ordem & Teste & Status esperado \\\midrule
01--03 & Saúde; listagem sem chave; listagem autenticada. & 200, 401, 200.\\
04--06 & Criar; ler o ID; repetir o nome. & 201, 200, 409.\\
07--08 & PATCH parcial; PUT completo. & 200, 200.\\
09--11 & Dados incompletos; formato errado; método errado. & 400, 415, 405.\\
12--13 & Excluir; ler o excluído. & 204, 404.\\\bottomrule
\end{tabularx}
Execute os itens em ordem, manualmente ou em \textit{Run Collection}, conforme os limites do plano. O POST usa um nome com instante atual para reduzir colisões entre execuções e salva o ID real em \cmd{equipamentoId}. Não rode duas execuções concorrentes no mesmo ambiente, pois elas compartilham essa variável.

\sub{Entenda um teste de resposta}
Na aba \textit{Scripts → Post-response} ou \textit{Tests}, conforme sua versão:
\begin{lstlisting}
pm.test("Criado", () => pm.response.to.have.status(201));
const dados = pm.response.json();
pm.environment.set("equipamentoId", String(dados.id));
\end{lstlisting}
Esse script é executado depois da resposta. Uma asserção aprovada verifica algo delimitado, como status ou campo. Não comprova que toda a infraestrutura é segura ou que todos os clientes funcionarão.

\sub{O caso 204 exige tratamento próprio}
\begin{lstlisting}
pm.test("Excluído sem corpo", () => {
  pm.response.to.have.status(204);
  pm.expect(pm.response.text()).to.eql("");
});
\end{lstlisting}
Não execute \cmd{pm.response.json()} indiscriminadamente para 204: não há JSON para interpretar. A interface web enfrenta o mesmo problema se chamar \cmd{response.json()} sempre.

\caixa{Verificação 5}{Entregue o resultado da execução, com os 13 itens e status esperados. Se um teste falhar, registre qual pré-condição não foi atendida. Um teste que espera 401 e recebe 401 deve passar; o objetivo não é obter 200 em todas as chamadas.}

\pagina{APLICAÇÃO WEB / CLIENTE PRONTO}{A página consome a mesma API}
A interface já foi publicada pelo script em \cmd{/painel.html}. Não há etapa de programação nem build. O código JavaScript do navegador envia requisições HTTP para a mesma API que você explorou no Postman.

\begin{enumerate}
\item Abra \cmd{https://SEU_DNS_EC2/painel.html}, informe a chave de g01 e liste os cadastros.
\item Abra DevTools → Network/Rede. Limpe a lista e crie um cadastro pela página. Localize POST \cmd{/api/equipamentos}: observe JSON, \cmd{Content-Type}, status 201 e o ID retornado. Oculte a chave nas evidências.
\item Consulte esse mesmo ID pelo Postman com g01. Deve retornar os mesmos dados. A interface não possui um inventário separado.
\item Altere \cmd{ativo} pelo Postman com PATCH. Liste novamente na página e confira a mudança. Faça o inverso: altere na página e confira por GET no Postman.
\item Exclua pela página com confirmação. Localize DELETE e o 204 sem corpo na aba Network; consulte o ID no Postman e confira 404.
\end{enumerate}

\sub{Do clique à representação}
\begin{tabularx}{\linewidth}{@{}lYY@{}}\toprule
Ação na página & Requisição observada & Evidência na resposta \\\midrule
Conectar / listar & GET da coleção & Lista JSON e total, 200.\\
Cadastrar & POST da coleção & ID do novo recurso, 201.\\
Ler um item & GET com ID & Representação do recurso, 200.\\
Salvar edição completa & PUT com ID & Representação atualizada, 200.\\
Alternar ativo & PATCH com ID & Campo alterado, 200.\\
Excluir & DELETE com ID & Corpo vazio, 204.\\\bottomrule
\end{tabularx}

\caixa{Verificação 6}{Escolha uma operação e explique a sequência: ação do usuário; método/URL/corpo enviados; resposta da API; atualização da tela. Compare o cadastro pelos dois clientes. O que permite criar outro aplicativo para a mesma API sem reescrever o armazenamento?}

\textbf{Por que isso permite construir aplicações web?} A API oferece um contrato de recursos e representações. Clientes distintos utilizam esse contrato por HTTP; a persistência e as regras de acesso permanecem no servidor. A página transforma ações do usuário em requisições e respostas em elementos visuais. REST orienta essa organização; não exige uma interface específica.

\textbf{Aprofundamento opcional:} o kit preserva o prompt completo do v0 no anexo, para quem quiser criar outra interface depois. Isso não é requisito de execução nem de entrega desta prática.

\pagina{NAVEGADOR / ORIGENS}{Evitar dois problemas comuns de integração}
\sub{Mesma origem: o caminho recomendado}
Uma origem é a combinação de esquema, host e porta. Se a página está em \cmd{https://SEU_DNS_EC2} e chama \cmd{/api/equipamentos}, a requisição usa a mesma origem. O Nginx encaminha internamente para o Python. A porta 8000 não aparece na URL pública e não é aberta para os alunos.

\sub{Origem diferente: o caso da prévia do v0}
Uma página em outro domínio que faça \cmd{fetch} para a EC2 está em outra origem. O navegador exige uma política CORS compatível. Com JSON e \cmd{X-API-Key}, normalmente envia antes OPTIONS (preflight), consultando método e cabeçalhos permitidos. A aplicação já sabe responder, mas uma origem externa só será aceita se cadastrada em \cmd{LAB_CORS_ORIGINS}. O caminho desta prática usa a mesma origem e dispensa esse ajuste.

\textbf{A chave continua necessária.} CORS é uma política aplicada por navegadores; não substitui autenticação nem um firewall. O Postman Desktop e o \cmd{curl} não aplicam a mesma restrição de origem que uma página. Consequentemente, funcionar no Postman não prova que funcionará no navegador.
\fonte{https://developer.mozilla.org/en-US/docs/Web/HTTP/Guides/CORS}{MDN: CORS e preflight}

\sub{HTTPS chamando HTTP: conteúdo misto}
Se a página é HTTPS e tenta chamar uma API HTTP remota, o navegador pode bloquear a requisição como conteúdo misto. Liberar CORS não corrige isso. Use HTTPS na API ou sirva a página e a API pela mesma origem HTTPS no Nginx.
\fonte{https://developer.mozilla.org/en-US/docs/Web/Security/Defenses/Mixed_content}{MDN: conteúdo misto}

\textbf{Certificado autoassinado em outra origem:} uma prévia hospedada pode falhar por confiança TLS, por CORS ou por ambos. Para a aula, teste o site exportado e a API juntos na EC2. Se a prévia externa for indispensável, use certificado publicamente confiável e permita somente a origem exata da página, após inspecionar o cabeçalho \cmd{Origin}. A origem pode mudar entre prévias.

\caixa{Não são correções válidas}{Não usar \cmd{mode: 'no-cors'}: a resposta fica opaca para o JavaScript e os cabeçalhos necessários são limitados. Não desabilitar a segurança do navegador e não expor a porta 8000 para contornar o proxy.}

\sub{Teste de distinção}
A API funciona no Postman, mas a página mostra ``Failed to fetch''. Verifique primeiro Console, Network, URL, certificado, esquema e OPTIONS. Se o servidor responde 401 em JSON, a requisição chegou à aplicação: investigue a chave.

\pagina{DIAGNÓSTICO E ENTREGA}{Transformar observação em explicação}
\begin{tabularx}{\linewidth}{@{}p{3.3cm}YY@{}}\toprule
Sintoma & O que observar & Próximo teste \\\midrule
Timeout & Não há resposta HTTP. & Rede/SG/rota; comparar com saúde local.\\
Erro de certificado & Falha antes de HTTP. & Confiança, SAN, data e nome/IP.\\
502 & Nginx responde; upstream falha. & Testar serviço e saúde em \cmd{127.0.0.1:8000}.\\
401 & A aplicação rejeita a chave. & Ambiente selecionado e header correto.\\
400 / 415 & Corpo ou formato incorreto. & Conferir JSON, tipos e Content-Type.\\
404 & Caminho/ID/grupo. & Listar os seus itens e usar ID real.\\
409 & Nome duplicado no grupo. & Consultar antes de repetir criação.\\
Só navegador falha & Console, OPTIONS e TLS. & Investigar CORS e conteúdo misto.\\
Interface mostra erro após excluir & DELETE retornou 204? & Não fazer parse de corpo ausente.\\\bottomrule
\end{tabularx}

\sub{Persistência e diagnóstico na sua EC2}
Crie um item com g01, anote o ID e execute \cmd{sudo systemctl restart redes-api} por SSH. Repita GET: espere 200 e os mesmos dados. Se houver 502, use \cmd{sudo systemctl status redes-api --no-pager} e \cmd{sudo journalctl -u redes-api -n 30 --no-pager}. Compare com a saúde local em \cmd{http://127.0.0.1:8000/api/saude}. Exclua o item após registrar a evidência.

\sub{Entrega do grupo}
\begin{enumerate}
\item Tabela de operações com método, caminho, corpo resumido, status e efeito confirmado por GET.
\item Resultado dos 13 testes da coleção ou execução manual equivalente.
\item Comparação fundamentada entre PUT/PATCH, 401/404 e falha de transporte/erro HTTP.
\item Evidência de persistência após reiniciar a API e isolamento entre g01/g02, sem exposição de chaves.
\item URL e captura do Network mostrando uma operação da página pronta sobre o mesmo cadastro acessado no Postman; explique como a interface utiliza o contrato REST.
\end{enumerate}

\sub{Perguntas para discussão}
Por que POST não é sinônimo de ``enviar qualquer coisa''? Por que sucesso no ping não prova o CRUD? Por que a chave não pertence à URL? Por que o Nginx pode ver o JSON mesmo quando o aluno usa HTTPS? Que evidência distingue 502 de 401?

\textbf{Ao encerrar:} exclua somente os cadastros do seu grupo; apague a chave de ambientes compartilhados; feche a interface; remova certificados/exceções didáticas adicionados à estação. Salve as evidências e termine a EC2 da sua equipe quando não for mais utilizada; confira recursos remanescentes e custos. Não há API docente compartilhada neste roteiro.

\textbf{Escopo:} scripts, API, página pronta e coleção acompanham o repositório/ZIP. O guia docente descreve a preparação da distribuição e os critérios de avaliação; nenhuma etapa do aluno depende de ações do professor no servidor.

{\footnotesize Revisão do fluxo autônomo em 23/10/2026; referências técnicas de 16--23/10/2026. Os nomes de abas e a disponibilidade de recursos do Postman/v0 variam conforme versão e plano.}
\pagina{ANEXO OPCIONAL / FORA DO NÚCLEO}{Criar outra interface com v0}
Copie este texto ou use o arquivo editável correspondente do kit. Não acrescente chaves reais. O aplicativo gerado deverá cumprir o contrato do backend fornecido.
\lstinputlisting[basicstyle=\ttfamily\footnotesize]{../kit/prompts/02-crud-v0.txt}
\finalizar
